\let\negmedspace\undefined
\let\negthickspace\undefined
\documentclass[journal]{IEEEtran}
\usepackage[a4paper, margin=10mm, onecolumn]{geometry}
%\usepackage{lmodern} % Ensure lmodern is loaded for pdflatex
\usepackage{tfrupee} % Include tfrupee package
\setlength{\headheight}{1cm} % Set the height of the header box
\setlength{\headsep}{0mm}   % Set the distance between the header box and the top of the text
\usepackage{gvv-book}
\usepackage{gvv}
\usepackage{cite}
\usepackage{amsmath,amssymb,amsfonts,amsthm}
\usepackage{algorithmic}
\usepackage{graphicx}
\usepackage{float}
\usepackage{textcomp}
\usepackage{xcolor}
\usepackage{txfonts}
\usepackage{listings}
\usepackage{enumitem}
\usepackage{mathtools}
\usepackage{gensymb}
\usepackage{comment}
\usepackage[breaklinks=true]{hyperref}
\usepackage{tkz-euclide} 
\usepackage{listings}
% \usepackage{gvv}                                       
\def\inputGnumericTable{}                                
\usepackage[latin1]{inputenc}                               
\usepackage{color}                                       
\usepackage{array}                                       
\usepackage{longtable}                                   
\usepackage{calc}                                        
\usepackage{multirow}                                    
\usepackage{hhline}                                      
\usepackage{ifthen}                                      
\usepackage{lscape}
\usepackage{tikz}
\usetikzlibrary{patterns}
\begin{document}
\bibliographystyle{IEEEtran}
\vspace{3cm}
\title{4.13.58}
\author{ee25btech11054-S.Harsha \vec{V}ardhan Reddy}
\maketitle
% \maketitle
% \newpage
% \bigskip
{\let\newpage\relax\maketitle}
\renewcommand{\thefigure}{\theenumi}
\renewcommand{\thetable}{\theenumi}
\setlength{\intextsep}{10pt} % Space between text and floats
\numberwithin{equation}{enumi}
\numberwithin{figure}{enumi}
\renewcommand{\thetable}{\theenumi}
\textbf{Question 4.13.58}
Show that all chords of the curve $3x^2 - y^2 - 2x + 4y = 0$, which subtend a right angle at the origin, pass through a fixed point. Find the coordinates of the point.

\solution
The equation of the given curve is
$3x^2 - y^2 - 2x + 4y = 0 $\\
Let $\vec{x} = \myvec{x \\ y}$.\\
The equation of the curve can be written in the matrix form $\vec{x}^\top {\vec{V}}\vec{x} + 2\vec{u}^\top\vec{x} + f = 0$ as
\begin{align}
\vec{x}^\top \myvec{3 & 0 \\ 0 & -1} \vec{x} + 2\myvec{-1 \\ 2}^\top\vec{x} + 0 = 0.
\end{align}
Here, the matrix of the quadratic part is 
\begin{align}
\vec{V} = \myvec{3 & 0 \\ 0 & -1}
\end{align}, 
the vector for the linear part is 
\begin{align}
\vec{u} = \myvec{-1 \\ 2}
\end{align}
and the constant term is $f=0$.

Let the equation of a chord PQ of the curve be $lx + my = 1 $.\\
In vector form, this is:
\begin{align}
\vec{n}^\top\vec{x} = 1,\\
\text{where } \vec{n} = \myvec{l \\ m}
\end{align}
The combined equation of the pair of lines joining the origin to the points of intersection, P and Q, is obtained by homogenizing equation (0.1) using equation (0.4).
\begin{align}
\vec{x}^\top {\vec{V}}\vec{x} + 2(\vec{u}^\top\vec{x})(1) = 0 
\end{align}
Substituting $\vec{n}^\top\vec{x}$ for $1$:
\begin{align}
\vec{x}^\top {\vec{V}}\vec{x} + 2(\vec{u}^\top\vec{x})(\vec{n}^\top\vec{x}) = 0 
\end{align}
This can be written as a single quadratic form:
\begin{align}
\vec{x}^\top {\vec{V}}\vec{x} + \vec{x}^\top(\vec{u}\vec{n}^\top + \vec{n}\vec{u}^\top)\vec{x} = 0\\
\vec{x}^\top( {\vec{V}} + \vec{u}\vec{n}^\top + \vec{n}\vec{u}^\top)\vec{x} = 0
\end{align}
This is the equation of the pair of lines OP and OQ. Let 
\begin{align}
{W} =  {\vec{V}} + \vec{u}\vec{n}^\top + \vec{n}\vec{u}^\top.
\end{align}
The lines represented by the homogeneous quadratic equation $\vec{x}^\top {W}\vec{x} = 0$ are perpendicular if and only if the trace of the matrix $ {W}$ is zero.
$\text{trace}( {W}) = 0 $
Using the property 
\begin{align}
\text{trace}( {A}+ {B}) = \text{trace}( {A}) + \text{trace}( {B})
\text{trace}(\vec{a}\vec{b}^\top) = \vec{b}^\top\vec{a}\\
\text{trace}( {W}) = \text{trace}( {\vec{V}}) + \text{trace}(\vec{u}\vec{n}^\top) + \text{trace}(\vec{n}\vec{u}^\top) \\
\text{trace}( {W}) = \text{trace}( {\vec{V}}) + \vec{n}^\top\vec{u} + \vec{u}^\top\vec{n} = \text{trace}( {\vec{V}}) + 2\vec{u}^\top\vec{n} 
\end{align}

The condition for the chord to subtend a right angle at the origin is 
\begin{align}
\text{trace}( {\vec{V}}) + 2\vec{u}^\top\vec{n} = 0 
\end{align}
From the problem statement:
\begin{align}
 {\vec{V}} = \myvec{3 & 0 \\ 0 & -1}\\
\implies \text{trace}( {\vec{V}}) = 3 + (-1) = 2 \\
\vec{u} = \myvec{-1 \\ 2} \\
\vec{n} = \myvec{l \\ m} 
\end{align}
Substituting these values into equation (0.17)
\begin{align}
2 + 2\myvec{-1 & 2}\myvec{l \\ m} = 0\\
2 + 2(-l + 2m) = 0\\
1 - l + 2m = 0 \\
l - 2m = 1 
\end{align}
This is the condition that the coefficients of the chord $lx+my=1$ must satisfy.
We want to show that this line passes through a fixed point $(x_0, y_0)$. If it does, then the coordinates of this point must satisfy the equation of the line for any $l$ and $m$ that satisfy condition (0.23).
\begin{align}
lx_0 + my_0 = 1
\end{align}
From condition (0.14), we express $l$ in terms of $m$: $l = 1 + 2m$. Substituting this into equation (0.24):
\begin{align}
(1 + 2m)x_0 + my_0 = 1 \\
x_0 + 2mx_0 + my_0 = 1 \\
(x_0 - 1) + m(2x_0 + y_0) = 0
\end{align}
Since this equation must hold for all possible values of $m$ (which corresponds to all possible chords subtending a right angle), the coefficients of the powers of $m$ must be individually zero.
\begin{align}
x_0 - 1 = 0\\
\implies x_0 = 1 \\
2x_0 + y_0 = 0 \\
\implies 2(1) + y_0 = 0\\
\implies y_0 = -2 
\end{align}
Therefore, all such chords pass through the fixed point with coordinates $\myvec{1\\-2}$.

\begin{figure}[h]
    \centering
    \includegraphics[width=1.1\linewidth]{figs/Figure_2.png}
    \caption{}
    \label{fig:placeholder}
\end{figure}

\end{document}